\section{问题三：烘干时长}\label{sec:q3}

\subsection{判据与结果}

烘干要求药材各处含水率低于 $\SI{0.15}{kg/kg}$，即 $\Cmax(t)<0.15$。由图~\ref{fig:fields}(b)，含水率沿半径由中心向表面单调减小，最大值始终在中心。在问题二的同一条解上求 $\Cmax(t)=0.15$ 的时刻，得到烘干时长
\begin{equation}
\tstar=\SI{57.4740}{h}\quad(\SI{206906.50}{s})
\end{equation}
约 2.4 天，与题目所说「烘干一般持续 2--3 天」一致。表~\ref{tab:t5} 给出每隔 $\SI{6}{h}$ 的含水率。result3.xlsx 每隔 $\SI{60}{s}$ 输出一行，末行为 $\SI{206940}{s}$，此时中心含水率为 0.149990，四舍五入后显示为 0.1500。

\begin{table}[H]\centering\small
\caption{药材烘干过程的水分浓度（\si{kg/kg}）}\label{tab:t5}
\begin{tabular}{cccccc}
\toprule
 & \multicolumn{5}{c}{到药材中心的距离/cm} \\
\cmidrule(lr){2-6}
时间/h & 0 & 0.5 & 1 & 1.5 & 2 \\
\midrule
6 & 1.0170 & 0.9881 & 0.9015 & 0.7550 & 0.5328 \\
12 & 0.4561 & 0.4432 & 0.4031 & 0.3297 & 0.1642 \\
18 & 0.2988 & 0.2916 & 0.2687 & 0.2253 & 0.0845 \\
24 & 0.2378 & 0.2327 & 0.2167 & 0.1855 & 0.0665 \\
30 & 0.2058 & 0.2018 & 0.1892 & 0.1641 & 0.0598 \\
36 & 0.1859 & 0.1825 & 0.1718 & 0.1504 & 0.0566 \\
42 & 0.1721 & 0.1691 & 0.1597 & 0.1407 & 0.0548 \\
48 & 0.1618 & 0.1592 & 0.1507 & 0.1334 & 0.0537 \\
54 & 0.1539 & 0.1514 & 0.1436 & 0.1277 & 0.0530 \\
57.4740（烘干结束） & 0.1500 & 0.1477 & 0.1402 & 0.1249 & 0.0526 \\
\bottomrule
\end{tabular}

\end{table}

\subsection{结果分析}

\parahead{中心决定烘干时长} 烘干结束时，表面含水率为 0.0526，接近烘房含水浓度，平均含水率为 0.1243，只有中心刚好降到 0.15。表面在 $\SI{12.6}{h}$、平均含水率在 $\SI{36.0}{h}$ 就已低于 0.15，若以表面或平均含水率判断，会把烘干时长分别低估 78\% 和 37\%。

\parahead{后期干燥很慢} 中心含水率在前 $\SI{6}{h}$ 由 2.55 降到 1.017，而从 $\SI{48}{h}$ 到烘干结束只下降 0.012。最大含水率降到 0.20、0.18、0.16 所需的时间分别为 $\SI{31.51}{h}$、$\SI{38.32}{h}$、$\SI{49.26}{h}$，从 0.16 降到 0.15 还要 $\SI{8.2}{h}$，占烘干时长的 14\%。烘干标准的微小变化会显著改变烘干时长，这是由式~\eqref{eq:twoterm} 中慢速项的长时间常数决定的。
